A otimização do descarregamento de tarefas em redes veiculares representa um desafio crítico para viabilizar sistemas de transporte inteligentes, exigindo decisões rápidas que equilibrem latência e consumo de energia. O avanço das arquiteturas de Computação de Borda Veicular (\acrotip{VEC}) transferiu a capacidade de processamento para mais perto dos usuários, mas a alta mobilidade e a variação da carga de rede tornam as heurísticas estáticas ineficientes. Por conta disso, a literatura recente tem transitado de métodos matemáticos tradicionais para abordagens baseadas em Aprendizado por Reforço Profundo (\acrotip{DRL}), buscando políticas de alocação que se adaptem dinamicamente ao ambiente.

O trabalho apresentado em~\cite{vieira2025foff} desenvolveu o framework \acrotip{FOFF} fundamentado em Tomada de Decisão Multicritério (\acrotip{MCDM}). A solução aplicou a técnica Fuzzy-TOPSIS para ranquear o servidor de borda (\acrotip{VEC}) e veículos vizinhos com base no tempo de resposta e consumo de energia. As simulações demonstraram que a abordagem obteve economias significativas de energia sem a necessidade de uma fase de treinamento prévio. Contudo, modelos baseados exclusivamente em regras lógicas ou avaliações heurísticas falham diante de cenários veiculares não estacionários, visto que não conseguem antecipar flutuações nas dinâmicas de fila ou aprender com o histórico de desempenho.

O estudo conduzido por~\cite{tan2018} explorou o aprendizado de máquina para lidar com a mobilidade dos usuários em redes veiculares. A pesquisa propôs um framework baseado em Redes Q Profundas (\acrotip{DQN}) para decidir a alocação em unidades de beira de estrada e gerenciar o descarregamento de tarefas. Os resultados indicaram uma redução efetiva nos custos de comunicação e cálculo. A principal limitação dessa abordagem reside em sua natureza centralizada, exigindo o controle simultâneo de todos os veículos da rede e provocando um aumento exponencial na complexidade computacional à medida que a densidade veicular cresce.

O trabalho de~\cite{pang2024} desenvolveu um algoritmo de descarregamento multitarefa utilizando a arquitetura Dueling Double Deep Q-Network (\acrotip{D3QN}). O método focou na minimização do tempo de resposta e do consumo energético em ambientes de borda altamente requisitados. As avaliações mostraram que o modelo reduziu a latência do sistema de forma mais eficaz que abordagens aleatórias tradicionais. No entanto, por ser baseado em espaços de ação estritamente discretos, o algoritmo apresenta dificuldades de escalabilidade quando aplicado a problemas de alocação contínua de recursos, limitando a precisão da distribuição de poder computacional.

A pesquisa apresentada por~\cite{fofana2024} propôs uma solução de descarregamento baseada em um jogo de Stackelberg de quatro estágios integrado a um modelo de aprendizado por reforço. A estratégia avalia o estado do veículo em tempo real para otimizar o tempo e a despesa do descarregamento de forma coordenada. Os experimentos confirmaram que a solução alcançou um equilíbrio na utilidade do sistema e melhorou a probabilidade de sucesso das tarefas. Apesar dos ganhos, a dependência de processos de negociação complexos entre múltiplos agentes introduz uma sobrecarga de sinalização que degrada o desempenho em cenários de tráfego denso e comunicação intermitente.

O estudo de~\cite{shi2023} buscou resolver o problema de alocação de recursos contínuos aplicando o algoritmo Ator-Crítico Determinístico Profundo (\acrotip{DDPG}). A abordagem utilizou a estrutura ator-crítico para reduzir o custo geral do sistema em um ambiente colaborativo de nuvem e borda veicular. O algoritmo convergiu rapidamente e reduziu os custos operacionais em comparação com métodos baseados apenas em valor. A falha dessa modelagem é a ausência de consideração sobre a dependência estrutural entre as tarefas e a utilização de preenchimento com zeros na observação do estado para lidar com o número dinâmico de veículos, obrigando a rede neural a processar artefatos espúrios e ruídos.

O trabalho desenvolvido por~\cite{gao2022} implementou um esquema adaptativo usando aprendizado por reforço multiagente em computação em névoa veicular. A solução distribuiu a tomada de decisão diretamente entre os veículos para reduzir a latência e o consumo de energia de forma descentralizada. As simulações demonstraram superioridade na capacidade de adaptação local em comparação com arquiteturas guiadas por agentes únicos. A limitação crítica dessa arquitetura é o desafio da sincronização e a dificuldade de garantir uma convergência estável quando múltiplos veículos alteram o ambiente simultaneamente, gerando instabilidade na política global de roteamento.

O estudo conduzido por~\cite{zheng2023} integrou a tecnologia de gêmeos digitais a redes de veículos para habilitar o descarregamento proativo assistido por inteligência artificial. O trabalho empregou o algoritmo Asynchronous Advantage Actor-Critic (\acrotip{A3C}) para otimizar as decisões baseando-se em uma réplica virtual contínua do ambiente físico veicular. A pesquisa relatou uma convergência rápida da rede e a diminuição considerável do custo do sistema. A restrição fundamental desse modelo é a altíssima sobrecarga computacional exigida para manter o gêmeo digital sincronizado em tempo real, fator que frequentemente inviabiliza sua execução fluida em dispositivos móveis.

A pesquisa de~\cite{nie2023} direcionou o aprendizado por reforço para tarefas de visão computacional e propôs um algoritmo focado na detecção de objetos em veículos autônomos. O modelo combinou \acrotip{DRL} com linearização por partes para maximizar especificamente uma função de utilidade baseada em tempo de execução. Os resultados apontaram melhorias expressivas no tempo de resposta para o processamento de imagens críticas. Contudo, a modelagem foi construída de forma monolítica para um tipo específico de aplicação, falhando em considerar as prioridades de diferentes tarefas concorrentes e negligenciando os problemas de congestionamento causados pela densidade veicular.

O estudo apresentado por~\cite{liu2024} focou no descarregamento de tarefas complexas modeladas como grafos direcionados acíclicos utilizando o algoritmo Ator-Crítico Suave (\acrotip{SAC}). O framework buscou equilibrar ativamente o atraso e o gasto de energia para aplicações com alta dependência estrutural entre seus subcomponentes. A avaliação atestou que a maximização da entropia estocástica evitou a convergência para ótimos locais e melhorou a utilidade de execução. A restrição principal desse modelo para o contexto automotivo é a premissa de um ambiente perfeitamente estático, ignorando completamente as taxas de desconexão e a variabilidade do canal de comunicação.

Uma abordagem híbrida foi desenvolvida por~\cite{wu2024}, que propôs o esquema \acrotip{TOLB} para descarregamento e balanceamento de carga integrando o Twin Delayed Deep Deterministic Policy Gradient (\acrotip{TD3}) com a Technique for Order of Preference by Similarity to Ideal Solution (\acrotip{TOPSIS}). A metodologia utilizou o algoritmo de aprendizado para decisões de proporção de particionamento e o cálculo multicritério para selecionar o servidor alvo. Os experimentos comprovaram a redução do atraso de processamento e a diminuição do desequilíbrio de carga na infraestrutura de borda. A limitação dessa abordagem reside no uso da filtragem multicritério apenas como um elemento de pós-processamento tático, mantendo o espaço de observação bruto e engessado na entrada do agente e prejudicando a generalização espacial em redes densas.

A arquitetura \acrotip{SAC}-\acrotip{TOPSIS} proposta nesta etapa intermediária difere substancialmente dessas abordagens, uma vez que consolida o método \acrotip{TOPSIS} internamente na borda de observação atuando estritamente como um módulo de engenharia de estado. Esse módulo avalia o servidor de borda (\acrotip{VEC}) e veículos vizinhos com base no Tempo de Conclusão de Tarefa (\acrotip{JCT}), consumo energético e distância ao nó de computação para identificar o melhor candidato ao descarregamento. A saída computada condensa essas características em uma codificação de estado contínua estritamente compacta como um vetor 6D perfeitamente invariante à escala, eliminando a dependência do preenchimento com zeros. A avaliação experimental comparou cinco famílias de \acrotip{DRL} (\acrotip{SAC}, Otimização de Política Proximal (\acrotip{PPO}), \acrotip{TD3}, \acrotip{DQN} e Advantage Actor-Critic (\acrotip{A2C})) e sete heurísticas de referência (\acrotip{VEC}, Random, Round-Robin, Gini, Jain, Greedy-Fair e Greedy-Unfair), atestando que o emparelhamento desse espaço estruturado produz ganhos consistentes. Isso valida a tese central de que a engenharia composicional do estado elaborada via \acrotip{TOPSIS} atua como o principal motor impulsionador da capacidade de generalização apresentada pelo agente inteligente em ambientes de tráfego variável. Em resumo, os principais aspectos dos trabalhos relacionados são apresentados na Tabela~\ref{tab:related_works}.

\begin{table*}[tpb]
\centering
\caption{Comparação das Soluções de Descarregamento -- Uma avaliação comparativa das estratégias de descarregamento propostas em estudos anteriores.}
\label{tab:related_works}
\renewcommand{\arraystretch}{1.3}
\begin{tabularx}{\textwidth}{c X X X X}
\hline
Ref. & Estratégia & Tecnologia de Comunicação & Conquista & Limitação \\
\hline
\cite{vieira2025foff} & Fuzzy-\acrotip{TOPSIS} & \acrotip{VEC} & Economias significativas de energia sem fase de treinamento prévio & Falha em adaptar-se a cenários veiculares não estacionários \\
\cite{tan2018} & \acrotip{DQN} & Redes Veiculares & Redução efetiva nos custos de comunicação e cálculo & Aumento exponencial na complexidade devido à natureza centralizada \\
\cite{pang2024} & \acrotip{D3QN} & \acrotip{IoV} e \acrotip{MEC} & Redução da latência do sistema de forma mais eficaz & Dificuldade de escalabilidade em problemas de alocação contínua \\
\cite{fofana2024} & \acrotip{DRL} + Jogo de Stackelberg & Redes Veiculares & Equilíbrio na utilidade do sistema e sucesso das tarefas & Alta sobrecarga de sinalização em cenários de tráfego denso \\
\cite{shi2023} & \acrotip{DDPG} & Nuvem e Borda Veicular & Convergência rápida e redução de custos operacionais & Utilização de preenchimento com zeros e processamento de ruídos \\
\cite{gao2022} & \acrotip{MARL} & Névoa Veicular & Superioridade na capacidade de adaptação local descentralizada & Desafio de sincronização e convergência instável na rede global \\
\cite{zheng2023} & \acrotip{A3C} + Gêmeos Digitais & \acrotip{IoV} & Convergência rápida da rede e diminuição do custo do sistema & Altíssima sobrecarga computacional exigida para sincronização \\
\cite{nie2023} & \acrotip{DRL} + Linearização por partes & Veículos Autônomos & Melhorias expressivas no tempo de resposta para processamento de imagens & Negligencia problemas de congestionamento e prioridades de tarefas \\
\cite{liu2024} & \acrotip{SAC} & Nuvem Veicular & Evita a convergência para ótimos locais e melhora a utilidade & Premissa de ambiente estático e ignorância da variabilidade do canal \\
\cite{wu2024} & \acrotip{TD3} + \acrotip{TOPSIS} & \acrotip{VEC} & Redução do atraso de processamento e do desequilíbrio de carga & Filtragem multicritério utilizada apenas como pós-processamento tático \\
\hline
\end{tabularx}
\end{table*}